\documentclass{article}

\usepackage[T1]{fontenc}
\usepackage{xcolor}
\usepackage{microtype}
\usepackage[colorlinks=true,linkcolor=blue,citecolor=teal]{hyperref}

\title{A Small Study of Reproducible Typesetting}
\author{Ada Example and Ben Reviewer}
\date{June 2026}

\begin{document}

\maketitle

\begin{abstract}
This note demonstrates a compact LaTeX workflow for comparing manuscript
versions. The revised version argues that reproducible typesetting depends on
declaring source files, generated bibliographies, and comparison artifacts as
explicit build dependencies.
\end{abstract}

\section{Motivation}

Scientific documents often combine prose, figures, tables, and bibliographic
records. When these ingredients are revised independently, reviewers need a
clear view of what changed. Tools such as \texttt{latexdiff} are useful because
they mark insertions and deletions directly in the compiled document.

The baseline workflow in this example is intentionally simple. Authors edit a
single \texttt{.tex} file, keep a generated \texttt{.bbl} file next to it, and
compile the document with \texttt{latexmk}. Classic references on literate
typesetting, LaTeX practice, and package maintenance motivate this approach
\cite{knuth1984,lamport1994,mittelbach2004}.

\section{Method}

The document is built from deterministic inputs. The bibliography is stored as
a checked-in \texttt{.bbl} file so that the example does not require running
BibTeX or Biber. This makes the comparison portable across small test
environments and keeps the diff focused on document content rather than tool
availability.

\section{Build rule}

The Makefile treats each logical document name as a target. For example,
\texttt{make old}, \texttt{make new}, and \texttt{make diff} all compile a
matching \texttt{.tex} file and count the same source with \texttt{texcount}.
The diff target first generates both \texttt{diff.tex} and \texttt{diff.bbl}.

\section{Conclusion}

The revised workflow is still small, but it makes the bibliography dependency
visible. That is the key detail needed for a reliable comparison target.

\input{\jobname.bbl}

\end{document}
